% ============================================================================
% Sección h) — Implementación en Python-Gurobi
% Lo que está entre [corchetes] se completa con la hoja "Resumen" de la corrida final.
% ============================================================================

\subsection{Implementación en Python-Gurobi}

El modelo se implementó en Python con las librerías \texttt{gurobipy} y \texttt{pandas}, separado en tres archivos: \texttt{datos.py} lee la instancia desde la carpeta \textit{Instancias Excel}, donde cada hoja corresponde a un conjunto o parámetro del modelo; \texttt{main.py} construye y resuelve el modelo, con las restricciones numeradas igual que en este informe; y \texttt{exportar.py} escribe los resultados en la carpeta \textit{Análisis de Excels}. Todas las rutas son relativas, por lo que el código se ejecuta de inicio a fin en cualquier computador.

La resolución se hace en dos etapas. En la primera se maximiza el material retirado, con un límite de 25 minutos. En la segunda se agrega la condición de no perder ese material y se minimiza el costo total, con un límite de 5 minutos. Esta segunda etapa evita que, cuando el presupuesto no es una restricción activa, la solución incluya contenedores o recorridos que no aportan material recuperado. La semilla y el número de hilos de Gurobi se fijan para que los resultados se puedan replicar exactamente.

Se resuelven dos escenarios con la misma instancia: el \textit{óptimo}, en que el modelo decide la ubicación de cada punto limpio, y el de \textit{ubicaciones actuales}, en que los siete puntos quedan donde están hoy y que sirve como caso base. Para reducir el tamaño del modelo, las variables de elección $q_{inm}$ solo se crean para los pares sector--sitio dentro de la distancia máxima, ya que fuera de ese radio el atractivo es nulo. Además, la restricción de presupuesto se expresa en millones de pesos para evitar problemas numéricos.

La Tabla~\ref{tab:tamano} resume el tamaño del modelo y el desempeño de Gurobi en la instancia del caso de estudio.

\begin{table}[h]
    \centering
    \caption{Tamaño del modelo y resultado de la resolución por escenario.}
    \label{tab:tamano}
    \begin{tabular}{lcc}
        \hline
        & Óptimo & Ubicaciones actuales \\
        \hline
        Variables totales            & [completar] & [completar] \\
        Variables binarias           & [completar] & [completar] \\
        Variables enteras            & [completar] & [completar] \\
        Restricciones                & [completar] & [completar] \\
        Tiempo de resolución [s]     & [completar] & [completar] \\
        Gap de optimalidad [\%]      & [completar] & [completar] \\
        \hline
    \end{tabular}
\end{table}
